\documentclass[11pt]{article}

\usepackage[a4paper,margin=2.4cm]{geometry}
\usepackage{amsmath,amssymb}
\usepackage[T1]{fontenc}
\usepackage{enumitem}
\usepackage{xcolor}
\usepackage[colorlinks=true,linkcolor=blue!45!black]{hyperref}

\newcommand{\eps}{\varepsilon}
\newcommand{\dd}{\mathrm{d}}
\newcommand{\code}{\smallskip\par\noindent\textbf{Computational task.}\ }
\newcommand{\deliverable}{\smallskip\par\noindent\textbf{Repo deliverable.}\ }
\newcommand{\challenge}{\smallskip\par\noindent\textbf{Challenge (optional).}\ }
\newcommand{\tut}[2]{\section*{Tutorial #1\quad{\normalsize\mdseries(#2)}}\addcontentsline{toc}{section}{Tutorial #1 --- #2}}

\setlist[enumerate]{itemsep=3pt,topsep=4pt}
\setlist[itemize]{itemsep=2pt,topsep=3pt}

\title{\vspace{-1.5em}\textbf{Mathematical Modelling --- Tutorials}\\[3pt]
\large Weeks 1--6}
\author{}
\date{}

\begin{document}
\maketitle
\vspace{-2em}

\noindent\rule{\textwidth}{0.4pt}
\begin{small}
\textbf{How these tutorials work.} Each tutorial has pen-and-paper parts and a \emph{computational task} that you carry out in whatever language you prefer (MATLAB is recommended; Python or Julia are equally fine). The point of the code is never the language --- it is that a classmate can regenerate your figure from your folder. Each week you add a folder \texttt{week-NN/<username>/} to the course tutorials repository, containing a half-page \texttt{report.md}, the figure(s) requested, the code that produced them, and a short \texttt{README} giving the command to regenerate them. You submit this as a pull request and review one classmate's in return (see the repository's \texttt{CONTRIBUTING.md}). Parts marked \emph{Challenge} are optional, for the motivated or when time allows. Chapter and appendix references are to the course notes. Every \texttt{report.md} should also close with one sentence of \emph{interpretation}: what does your result mean for the real system you modelled, and not merely as a piece of mathematics?
\end{small}
\noindent\rule{\textwidth}{0.4pt}

% ======================================================================
\tut{1}{Week 1 --- Dimensional analysis in practice}

\emph{Notes: Chapter on dimensional analysis.} This week also serves as your onboarding to the course repository.

\begin{enumerate}
\item Rederive, by dimensional analysis alone, the period of a simple pendulum $T=\sqrt{L/g}\,\Phi(\theta_0)$, stating clearly what the method fixes and what it leaves in the unknown function $\Phi$.
\item Pick one further system --- the frequency of a mass on a spring, the speed of deep-water waves, or the draining time of a tank --- and carry out the same analysis: list the variables, form the dimensionless groups, and state precisely what is predicted and what is left undetermined.
\item A rigid-body (compound) pendulum --- an extended rigid object swinging about a pivot --- has small-amplitude period $T=2\pi\sqrt{L_{\mathrm{eff}}/g}$, where $g$ is gravity and $L_{\mathrm{eff}}=I/(mR)$ is its \emph{effective length}, built from the body's mass $m$, its moment of inertia $I$ about the pivot, and the pivot-to-centre-of-mass distance $R$ (see the notes, Chapter 6). Without computing any moment of inertia, explain why dimensional analysis already guarantees the $\sqrt{L_{\mathrm{eff}}/g}$ dependence.
\end{enumerate}

\code Taylor's blast-wave argument predicts $R=C\,(E t^2/\rho)^{1/5}$ for the radius $R$ of a blast at time $t$. Generate a synthetic ``photographic'' data set: fix $\rho=1.2$ and take the constant $C=1$ (its theoretical value), choose an energy $E$, compute $R(t)$ at a dozen times, and add a few percent of random noise to mimic measurement error. Now \emph{pretend you do not know $E$}: fit a straight line to $\log R$ versus $\log t$, check that the slope is close to $2/5$, and recover $E$ from the intercept (with $C=1$ and $\rho$ known, the intercept gives $E/\rho$). Report how close your estimate is to the true $E$.

\deliverable Your fitting code, a log--log plot showing the data and the fitted line, and a \texttt{report.md} stating the recovered slope and energy --- closing with what this shows: that a handful of photographs, plus dimensional analysis alone, suffice to weigh an explosion. This is your first pull request --- get the workflow working now, while the mathematics is easy.

% ======================================================================
\tut{2}{Week 2 --- Nondimensionalisation and the phase plane}

\emph{Notes: Chapter on nondimensionalisation; Appendix on linear stability.}

\begin{enumerate}
\item The pendulum can be written $\ddot\theta+\omega^2\sin\theta=0$, where $\omega^2=g/L$ is the squared natural frequency --- this is simply the notes' equation $L\ddot\theta+g\sin\theta=0$ divided by $L$. Show that rescaling time by $1/\omega$ turns it into $\theta''+\sin\theta=0$, with $\omega$ gone entirely. Conclude that setting $\omega=1$ loses no generality --- any $\omega$ can be scaled away --- which is why we may always do it. (This answers the standing question, ``why may we set $\omega=1$?'')
\item Write the nondimensional damped pendulum $\theta''+c\,\theta'+\sin\theta=0$, in which $c\ge0$ measures the damping, as the first-order system $\dot x=y,\ \dot y=-\sin x-c\,y$. Find all equilibria. Using the Jacobian and the trace--determinant classification (Appendix A), determine the type and stability of the downward equilibrium $(0,0)$ for $c=0$, for $0<c<2$, and for $c>2$, and of the upright equilibrium $(\pi,0)$. Compute the eigenvectors at the upright saddle.
\end{enumerate}

\code Plot the phase portrait (direction field plus a few trajectories) of the damped pendulum for $c=0$, $c=0.3$, and $c=3$. (MATLAB users may find the ready-made \texttt{dfield}/\texttt{pplane} tools convenient; in another language, a short script pairing a direction-field plot --- a \texttt{quiver} or \texttt{streamplot} --- with an ODE integrator does the same job.) Confirm visually the centre, the stable spiral, and the stable node at the downward equilibrium, and the saddle at the upright one.

\challenge Add a control torque to make the \emph{upright} position stable: with $F=a\,\phi+b\,\dot\phi$ (where $\phi=\theta-\pi$ measures displacement from upright), the linearised motion becomes $\ddot\phi+p\,\dot\phi+q\,\phi=0$. Find $p$ and $q$ in terms of $a,b,c,m,L,g$, and shade the region of the $(a,b)$-plane where the upright position is stable (i.e.\ $p>0$ and $q>0$). Interpret: what must the controller supply?

\deliverable The phase-portrait figure(s), and a \texttt{report.md} listing the equilibria with their types --- and saying what each means physically (the bob hanging straight down is stable; balanced upright it is not, which is why the challenge needs active control). If you did the challenge, include the $(a,b)$ stability region.

% ======================================================================
\tut{3}{Week 3 --- Amplitude and period of a nonlinear oscillator}

\emph{Notes: Chapter on regular perturbation --- Poincar\'e--Lindstedt and the energy method.}

For the undamped pendulum $\ddot\theta+\sin\theta=0$ released from rest at amplitude $\theta_0$:
\begin{enumerate}
\item Using energy conservation, show that $\dot\theta^2=2(\cos\theta-\cos\theta_0)$, and hence that the period is the quadrature
\[
T(\theta_0)=4\int_0^{\theta_0}\frac{\dd\theta}{\sqrt{2(\cos\theta-\cos\theta_0)}} .
\]
\item Recall the Poincar\'e--Lindstedt prediction $T\approx 2\pi\bigl(1+\tfrac{1}{16}\theta_0^2\bigr)$, and the small-amplitude (linear) value $T=2\pi$.
\end{enumerate}

\code Evaluate the exact period $T(\theta_0)$ numerically for $\theta_0$ from small up to near $\pi$ (watch the integrable singularity at the upper limit). On one figure plot three curves: the exact $T(\theta_0)$, the linear value $2\pi$, and the Poincar\'e--Lindstedt approximation. Identify the amplitude beyond which the perturbative curve visibly departs from the exact one --- this is the function $\Phi(\theta_0)$ of Tutorial~1, now computed.

\deliverable The period-versus-amplitude figure and a \texttt{report.md} commenting on the range of validity of the small-amplitude and Poincar\'e--Lindstedt models.

% ======================================================================
\tut{4}{Week 4 --- Slow--fast systems and relaxation oscillations}

\emph{Notes: Chapter on multiple scales and slow--fast reduction.}

Consider the Van der Pol oscillator $\ddot v-\mu(1-v^2)\dot v+v=0$.
\begin{enumerate}
\item Put it in Li\'enard slow--fast form and identify the critical (slow) manifold $y=\tfrac13 v^3-v$ and its folds at $v=\pm1$. Which branches are attracting?
\item Predict, from the shape of the manifold, the qualitative motion for large $\mu$ (slow crawls, fast jumps).
\end{enumerate}

\code Integrate the Van der Pol equation numerically for $\mu=0.1,\ 1,\ 5,\ 20$. For each, plot $v(t)$ and the phase-plane trajectory, and comment on the transition from near-sinusoidal to relaxation behaviour. Overlay the critical manifold on the phase plane for the large-$\mu$ case, and estimate the period as a function of $\mu$; compare with the large-$\mu$ asymptotic period $(3-2\ln2)\mu$.

\challenge Start the trajectory \emph{off} the critical manifold and show the fast initial transient that quasi-steady-state reduction misses; explain the failure using Tikhonov's condition from the notes.

\deliverable The waveform and phase-portrait figures, and a \texttt{report.md} on the sinusoidal-to-relaxation transition and your period estimates.

% ======================================================================
\tut{5}{Week 5 --- Building a model: traffic}

\emph{Notes: Chapter on balance laws and closures; bifurcations as predictions.}

The car-following model is $\dot h_i=v_{i+1}-v_i$, $\dot v_i=\alpha\bigl(V(h_i)-v_i\bigr)$, with the optimal-velocity closure
\[
V(h)=V_t\,\frac{[(h-h_{\min})/s]^k}{1+[(h-h_{\min})/s]^k}\quad(h>h_{\min}),\qquad V=0\ (h\le h_{\min}).
\]
\begin{enumerate}
\item Plot $V(h)$ for $h_{\min}=10\,$m, $V_t=100\,$km/h, and explore its shape: $k=1$ with $s=1$, $s<1$, $s>1$; then $s=1$ with $k=2,3,4$. Argue why $k=3,\,s=1$ makes a reasonable model of a ``normal'' driver.
\item On a ring road of length $L$ with $n$ cars, the uniform-flow state has headway $h^\ast=L/n$ and speed $V(h^\ast)$. Show that gridlock (all cars stopped) occurs when $L\le n\,h_{\min}$. Derive the flow rate $r(h)=V(h)/h$ and show its maximum satisfies $V'(h^\ast)=V(h^\ast)/h^\ast$.
\end{enumerate}

\code Simulate the ring-road system for, say, $n=20$ cars. Start close to uniform flow, add a small perturbation to one car, and integrate. Show that smooth flow persists when $V'(h^\ast)<\tfrac12\alpha$ but breaks into a travelling stop-and-go wave when $V'(h^\ast)>\tfrac12\alpha$ (a phantom jam). Separately, plot the fundamental diagram $r(h)$ and mark the maximum-flow headway.

\deliverable A space--time plot (car position or headway versus time) showing a jam forming, the fundamental-diagram figure, and a \texttt{report.md} relating the jam onset to the stability criterion.

% ======================================================================
\tut{6}{Week 6 --- Reading and improving a model (capstone launch)}

\emph{Notes: Chapter on reading and improving a model. This tutorial launches the half-assessment capstone; see the capstone brief for the full requirements.}

The capstone asks you to reconstruct a model, criticise it, and improve it so that it makes a new, checkable prediction. You may build it on \emph{either} of two models (as set out in the capstone brief):
\begin{itemize}
\item \textbf{traffic} --- the car-following model of Tutorial~5 (recommended, since you have already built and simulated it); or
\item \textbf{the Van der Pol oscillator} --- the circuit model of the notes, extended by forcing or coupling.
\end{itemize}
This week you take the first steps for your chosen model.
\begin{enumerate}
\item \emph{Reconstruct.} From the stated assumptions, rebuild the model in your own notation, and separate the exact balance law from the constitutive closure. (Traffic: headway bookkeeping versus the optimal-velocity function $V(h)$. Van der Pol: Kirchhoff's current law versus the nonlinear characteristic $f$, yielding $\ddot v-\mu(1-v^2)\dot v+v=0$.)
\item \emph{Critique, and choose an improvement.} Identify the weakest assumption and propose one improvement --- stated as a \emph{prediction}, a qualitative transition your improved model forecasts. (Traffic: a reaction-time delay, predicting a shifted jamming threshold; or a slow driver, predicting a bottleneck. Van der Pol: forcing, predicting frequency locking; or coupling, predicting an entrainment threshold.)
\end{enumerate}

\code Reproduce one result of the base model numerically --- the ring-road jam of Tutorial~5, or the relaxation waveform of the oscillator. Then take the first step of your improvement and produce a figure that begins to test its prediction. This is the seed of your capstone repository.

\deliverable Start your own capstone repository under the course organisation, containing the reproduced base result, a first improvement figure, and a \texttt{README} outlining your intended critique and improvement. (This grows into the half-assessment submission.)

% ======================================================================
\tut{$\star$}{Optional challenge --- Modelling on networks}

\emph{Not required; a self-contained topic for the motivated, or if time allows. It is deliberately discrete, in contrast to the continuous models of the course.}

Consider a transport network specified by a matrix of edge capacities $c_{ij}$, with a source node $s$ and a target node $t$ (your instructor will provide a specific network, or invent your own $7$-node example).
\begin{enumerate}
\item Draw the network from its capacity matrix.
\item Interpreting $c_{ij}$ as the capacity of the edge from $i$ to $j$, find the maximum flow that can be routed from $s$ to $t$, and identify the bottleneck (the minimum cut).
\item What happens to the maximum flow if every capacity is scaled by a positive constant $b$? If every capacity is squared?
\end{enumerate}

\code Implement, or use a library routine for, the max-flow (and, if you like, shortest-path) computation. Verify it by hand on a small graph first, then explore the scaling questions above numerically. A short discussion of the max-flow/min-cut duality makes a good report.

\deliverable Your code, a drawing of the network with the maximum flow marked on each edge, and a \texttt{report.md} answering the scaling questions.

\end{document}